\documentclass[a4paper]{scrartcl}
\usepackage[
    typ=ueb,
    fach=Informatik,
    lerngruppe=9g,
    loesungen=seite,
    module={Symbole,Bewertung},
]{schule}

% Dieses Dokument gehört zu den Beispiel des LaTeX Paketes Schule und ist von den Autoren
% des Pakets erstellt worden.
%
% Das Dokument steht unter der Lizenz: Creative Commons by-nc-sa Version 4.0
% http://creativecommons.org/licenses/by-nc-sa/4.0/deed.de
%
% Nach dieser Lizenz darf das Dokument beliebig kopiert und bearbeitet werden,
% sofern das Folgeprodukt wiederum unter gleichen Lizenzbedingungen vertrieben
% und auf die ursprünglichen Urheber verwiesen wird.
% Eine kommerzielle Nutzung ist ausdrücklich ausgeschlossen.

\author{Test Person}
\date{\today}
\title{Ein Übungsblatt}

\usepackage{blindtext}

\begin{document}

\section*{\Titel}
\blindtext

\section*{Übungen}
\begin{aufgabe}
    Eine Übung.
\end{aufgabe}
\begin{loesung}
    Erste Lösung
\end{loesung}

\begin{aufgabe}
    \setzeSymbol{\symHaken}
    Eine Übungsaufgabe mit zwei Teilaufgaben
    \begin{teilaufgaben}
        \teilaufgabe Erste Teilaufgabe.
        \teilaufgabe Zweite Teilaufgabe.
    \end{teilaufgaben}
\end{aufgabe}
\begin{loesung}
    Lösung für die zwei Teilaufgaben
\end{loesung}

\begin{aufgabe*}[points=10,bonus-points=5]
    Eine Zusatzübung.
    \begin{loesung*}
        Eine weitere Lösung.
    \end{loesung*}
    \begin{bearbeitungshinweis}
        Ein Hinweis für die Zusatzübung
    \end{bearbeitungshinweis}
\end{aufgabe*}

\end{document}
